\documentclass[11pt,a4paper]{article}
\usepackage[utf8]{inputenc}
\usepackage[italian]{babel}
\usepackage{geometry}
\usepackage{titlesec}
\usepackage{enumitem}
\usepackage{booktabs}
\usepackage{tabularx}
\usepackage{color}
\usepackage{colortbl}
\usepackage{tcolorbox}
\usepackage{fancyhdr}
\usepackage{graphicx}
\usepackage{amssymb}

\geometry{margin=1.5cm}

\definecolor{romgreen}{RGB}{0, 100, 0}
\definecolor{romlightgreen}{RGB}{0, 130, 0}
\definecolor{ferengilightgold}{RGB}{255, 215, 100}
\definecolor{ferengigold}{RGB}{204, 153, 0}

\titleformat{\section}
  {\normalfont\Large\bfseries\color{romgreen}}
  {}{0em}{}[\titlerule]

\titleformat{\subsection}
  {\normalfont\large\bfseries\color{romgreen}}
  {}{0em}{}

\pagestyle{fancy}
\fancyhf{}
\renewcommand{\headrulewidth}{1pt}
\renewcommand{\footrulewidth}{1pt}
\fancyhead[L]{RRW Thalassar}
\fancyhead[C]{Alleanza Romulano-Ferengi}
\fancyhead[R]{Scum \& Villainy}
\fancyfoot[C]{\thepage}

\begin{document}

\begin{center}
\Huge\textbf{\textcolor{romgreen}{RRW Thalassar (Classe Tal'Aura)}}\\[0.5cm]
\Large\textbf{Pannello Master Romulano-Ferengi}
\end{center}
\begin{tcolorbox}[colback=romlightgreen!20, colframe=romgreen, title=\textbf{Informazioni sulla Nave}]
\begin{tabular}{ll}
\textbf{Classe}: & Tal'Aura \\
\textbf{Equipaggio}: & 200-250 membri (alta % di personale scientifico) \\
\textbf{Velocità massima}: & Curvatura 9.7 \\
\textbf{Armamento}: & Disruptori di precisione, siluri al plasma \\
\textbf{Specializzazione}: & Archeologia, ingegneria inversa, occultamento \\
\end{tabular}
\end{tcolorbox}

\section{Quantum Data Points (QDP)}

\begin{tcolorbox}[colback=romlightgreen!10, colframe=romgreen]
\begin{tabularx}{\textwidth}{|X|X|X|}
\hline
\rowcolor{romgreen!20} \textbf{QDP Alpha} & \textbf{QDP Beta} & \textbf{QDP Gamma} \\
\hline
\large $\square\square$ \textbf{0} & \large $\square\square$ \textbf{2} & \large $\square\square$ \textbf{0} \\
\hline
\textbf{Conoscenza, Analisi,} & \textbf{Tecnologia, Potere,} & \textbf{Risorse, Materiali,} \\
\textbf{Scienza} & \textbf{Manipolazione} & \textbf{Estrazione} \\
\hline
\end{tabularx}
\end{tcolorbox}

\section{Dualità Romulano-Ferengi}

\begin{tcolorbox}[colback=romgreen!10, colframe=romgreen, title=\textbf{Approccio Romulano}]
\begin{itemize}[leftmargin=*]
\item Segretezza e manipolazione strategica
\item Pazienza e pianificazione a lungo termine
\item Interesse nel potere militare e strategico dell'artefatto
\item Lealtà all'Impero (frammentato) e alle proprie case nobili
\item Cautela e diffidenza verso le altre fazioni
\end{itemize}
\end{tcolorbox}

\begin{tcolorbox}[colback=ferengigold!20, colframe=ferengigold, title=\textbf{Approccio Ferengi}]
\begin{itemize}[leftmargin=*]
\item Opportunismo e ricerca del profitto immediato
\item Negoziazioni aggressive e "interpretative"
\item Interesse nel potenziale commerciale dell'artefatto
\item Lealtà al proprio portafoglio e alle Regole dell'Acquisizione
\item Pragmatismo e adattabilità nelle alleanze
\end{itemize}
\end{tcolorbox}

\section{Meccaniche del Gioco}

\begin{tcolorbox}[colback=romlightgreen!10, colframe=romgreen]
\begin{itemize}[leftmargin=*]
\item Focus sui QDP Beta (tecnologia, potere)
\item Enfatizzare il conflitto interno tra approcci romulani e ferengi
\item Opportunità per strade divergenti e strategie doppie
\item Tentativi di manipolazione e sfruttamento delle altre fazioni
\item Bilanciare obiettivi a lungo termine (romulani) e guadagni immediati (ferengi)
\end{itemize}
\end{tcolorbox}

\section{Enigmi della RRW Thalassar}

\begin{tcolorbox}[colback=romlightgreen!10, colframe=romgreen]
\begin{description}[leftmargin=*]
\item[La Risonanza dell'Occultamento] Il dispositivo di occultamento inizia a risuonare con l'artefatto, creando microscopiche "bolle temporali" attorno alla nave.
\item[Il Paradosso del Profitto] Cristalli iconiani formano visioni contraddittorie di possibili futuri per Ferengi e Romulani.
\item[Lo Specchio dei Desideri] Figure spettrali iconiane mostrano a ciascun membro dell'equipaggio visioni di ciò che desiderano maggiormente.
\item[L'Eco Quantico] Il nucleo a singolarità quantica vibra in risonanza con l'artefatto, producendo echi con ritardi precisi.
\end{description}
\end{tcolorbox}

\section{Capacità Uniche}

\begin{tcolorbox}[colback=romlightgreen!10, colframe=romgreen]
\begin{description}[leftmargin=*]
\item[Dispositivo di Occultamento] Permette di operare non rilevati dalle altre navi (costo: 2 QDP Beta)
\item[Infiltrazione Mascherata] Teletrasporto più difficile da rilevare (costo: 3 QDP)
\item[Scansioni Occultate] Possono analizzare altre navi senza essere rilevati
\item[Tecnologia di Singolarità] Sistemi energetici unici che possono interagire con l'artefatto iconiano
\end{description}
\end{tcolorbox}

\section{Fasi dell'Avventura}

\begin{tcolorbox}[colback=romlightgreen!10, colframe=romgreen]
\begin{enumerate}[leftmargin=*]
\item \textbf{Introduzione (15 minuti)}: Stabilire dinamiche interne e prime strategie
\item \textbf{Missioni Iniziali (60 minuti)}: Recupero e analisi di un frammento dell'artefatto
\item \textbf{Svolta (90 minuti)}: Possibilità di manipolare l'attivazione del portale
\item \textbf{Conclusione (75 minuti)}: Scelta strategica finale per ottenere vantaggio dall'artefatto
\end{enumerate}
\end{tcolorbox}

\section{Strategie di Doppio Gioco}

\begin{tcolorbox}[colback=romlightgreen!10, colframe=romgreen]
\begin{itemize}[leftmargin=*]
\item Apparire cooperativi mentre si perseguono agende nascoste
\item Condividere informazioni imprecise o incomplete con altre navi
\item Negoziare accordi separati con diverse fazioni
\item Sfruttare punti di debolezza delle altre navi (scientifica/militare/mineraria)
\item Infiltrare altre navi tramite teletrasporto quando possibile
\end{itemize}
\end{tcolorbox}

\section{Personaggi Chiave}

\begin{tcolorbox}[colback=romlightgreen!10, colframe=romgreen]
\begin{description}[leftmargin=*]
\item[Comandante Tr'Ellor] Ufficiale romulano, veterano diplomatico con legami al Tal Shiar
\item[DaiMon Brek] Leader ferengi, commerciante audace e abile negoziatore
\item[Sottocomandante Verat] Ufficiale scientifico romulano, esperto di tecnologie aliene
\item[Consulente commerciale Nog] Ferengi con esperienza nella Federazione, più moderato
\item[Agente Tal'Venna] Agente del Tal Shiar sotto copertura, con agenda segreta
\end{description}
\end{tcolorbox}

\section{Meccanica del Teletrasporto in Presenza di Interferenze Iconiane}

\begin{tcolorbox}[colback=romlightgreen!10, colframe=romgreen]
\begin{itemize}[leftmargin=*]
\item \textbf{Costo Base}: 3 QDP
\item \textbf{Capacità Unica}: "Infiltrazione Mascherata" - Più difficile da rilevare (3 QDP)
\item \textbf{Bonus}: Occultamento riduce la difficoltà del teletrasporto verso altre navi
\item \textbf{Risoluzione}: 2d6 + modificatori, dove 10+ è successo completo
\end{itemize}
\end{tcolorbox}

\section{Opzioni per la Fase Finale}

\begin{tcolorbox}[colback=romlightgreen!10, colframe=romgreen]
\begin{description}[leftmargin=*]
\item[Controllo Completo] (7 QDP Beta) Tentare di assumere il controllo esclusivo del portale
\item[Riprogrammazione Selettiva] (4 QDP Beta + 3 QDP Alpha) Configurare il portale per destinazioni strategiche
\item[Sabotaggio Controllato] (5 QDP Beta + 2 QDP Gamma) Danneggiare il sistema per impedire che altri possano usarlo
\end{description}
\end{tcolorbox}

\section{Suggerimenti di Gioco}

\begin{tcolorbox}[colback=romlightgreen!10, colframe=romgreen]
\begin{itemize}[leftmargin=*]
\item Crea tensione interna tra membri romulani e ferengi dell'equipaggio
\item Alterna momenti di effettiva collaborazione a momenti di sospetto reciproco
\item Permetti ai giocatori di scegliere quale approccio culturale privilegiare
\item Usa la tecnologia di occultamento come vantaggio strategico chiave
\item Offri opportunità per manipolare le altre fazioni e metterle l'una contro l'altra
\end{itemize}
\end{tcolorbox}

\section{Regole dell'Acquisizione Rilevanti}

\begin{tcolorbox}[colback=ferengigold!20, colframe=ferengigold]
\begin{itemize}[leftmargin=*]
\item \#1: Una volta che hai i loro soldi, non restituirli mai.
\item \#3: Non spendere mai più di quanto devi.
\item \#9: Opportunità più dimensioni uguali profitto.
\item \#34: La pace è buona per gli affari.
\item \#35: La guerra è buona per gli affari.
\item \#47: Non fidarti mai di un uomo che indossa un abito migliore del tuo.
\item \#194: È sempre meglio conoscere i tuoi clienti prima che conoscano te.
\end{itemize}
\end{tcolorbox}

\section{Comandi del Bot}

\begin{tcolorbox}[colback=romlightgreen!10, colframe=romgreen]
\begin{description}[leftmargin=*]
\item[/register] Registrazione come Master
\item[/status] Stato corrente dell'avventura
\item[/scan] Eseguire una scansione dell'artefatto
\item[/qdp\_status] Vedere lo stato dei tuoi QDP
\item[/add\_qdp <tipo> <quantità> <descrizione>] Aggiungere QDP (alpha, beta, gamma)
\item[/spend\_qdp <tipo> <quantità> <descrizione>] Spendere QDP
\item[/resolve\_enigma\_piece <tipo>] Risolvere parte di un enigma
\item[/teleport <tavolo> <descrizione>] Tentare un teletrasporto
\item[/cloak] Attivare/disattivare il dispositivo di occultamento
\item[/cloaking\_status] Controllare lo stato del dispositivo di occultamento
\item[/convergence\_roll <descrizione>] Eseguire un tiro di Convergenza nella fase finale
\item[/send\_message <tavolo> <messaggio>] Inviare un messaggio ad un altro tavolo
\end{description}
\end{tcolorbox}

\end{document}